\documentclass[11pt,a4paper]{article}
% Source of public/resume.pdf. Rebuild from this folder with:
%   pdflatex resume.tex   (then copy resume.pdf to ../public/resume.pdf)

\usepackage[utf8]{inputenc}
\usepackage[T1]{fontenc}
\usepackage[english]{babel}
\usepackage[a4paper,margin=1.7cm,top=1.0cm,bottom=1.1cm]{geometry}
\usepackage{mathpazo}
\usepackage{microtype}
\usepackage{xcolor}
\usepackage{enumitem}
\usepackage{tabularx}
\usepackage{array}
\usepackage{hyperref}
\usepackage{tikz}
\usepackage[most]{tcolorbox}
\usepackage{ragged2e}
\usepackage{pifont}
\usepackage{textcomp}

% ---------------------------------------------------------------- colours
\definecolor{MainBlue}{RGB}{18,55,105}
\definecolor{AccentBlue}{RGB}{40,110,175}
\definecolor{DarkGray}{RGB}{60,60,65}
\definecolor{LightGray}{RGB}{235,237,242}
\definecolor{ChipBg}{RGB}{225,232,243}

\hypersetup{
  colorlinks=true,
  urlcolor=AccentBlue,
  linkcolor=AccentBlue,
  pdfauthor={Elie Rouphael},
  pdftitle={Elie Rouphael - Curriculum Vitae},
  pdfsubject={Curriculum Vitae},
  pdfkeywords={system identification, model predictive control, learning, dynamical systems}
}

\pagestyle{empty}
\setlength{\parindent}{0pt}
\setlength{\JustifyingParindent}{0pt}
\setlength{\parskip}{0pt}
\setlist[itemize]{leftmargin=1.15em,itemsep=1pt,topsep=1pt,parsep=0pt,partopsep=0pt,
                  label=\textcolor{AccentBlue}{\small\textbullet}}
\renewcommand{\arraystretch}{1.12}
\newcolumntype{L}[1]{>{\raggedright\arraybackslash}p{#1}}
\newcolumntype{Y}{>{\raggedright\arraybackslash}X}

% ---------------------------------------------------------------- commands
\newcommand{\cvsection}[1]{%
  \vspace{0.6em}%
  \noindent{\large\bfseries\color{MainBlue}#1}\par
  \vspace{0.15em}{\color{AccentBlue}\hrule height 1.1pt}%
  \vspace{0.35em}}

\newcommand{\cvitem}[2]{%
  \textbf{\color{MainBlue}#1} & #2\\}

\newcommand{\datebadge}[1]{%
  \tikz[baseline=(d.base)]{%
    \node[fill=MainBlue,text=white,rounded corners=2pt,
          inner xsep=5pt,inner ysep=1.5pt,font=\footnotesize\bfseries] (d) {#1};
  }}

\newcommand{\role}[4]{%
  \begin{minipage}[t]{\dimexpr\textwidth-3.9cm\relax}
    \textbf{\color{MainBlue}#1}\\[0.5pt]
    \textit{\color{DarkGray}#3}
  \end{minipage}\hfill
  \begin{minipage}[t]{3.8cm}\raggedleft
    \datebadge{#2}
  \end{minipage}\\[2pt]
  #4\vspace{0.2em}}

\newcommand{\chip}[1]{%
  \tikz[baseline=(c.base)]{%
    \node[fill=ChipBg,text=MainBlue,rounded corners=2pt,
          inner xsep=4pt,inner ysep=1pt,font=\footnotesize\bfseries] (c) {#1};
  }}

\newcommand{\pubnum}[1]{%
  \tikz[baseline=(p.base)]{%
    \node[draw=AccentBlue,fill=white,text=AccentBlue,rounded corners=2pt,
          inner xsep=3pt,inner ysep=0.5pt,font=\footnotesize\bfseries,
          line width=0.5pt] (p) {#1};
  }}

\newcommand{\coursera}[2]{\href{https://www.coursera.org/account/accomplishments/verify/#2}{#1}}

\begin{document}

% ---------------------------------------------------------------- header
\begin{tcolorbox}[
  enhanced,
  colback=LightGray,
  colframe=LightGray,
  arc=3pt,
  boxrule=0pt,
  left=12pt,right=12pt,top=10pt,bottom=10pt,
  borderline west={3pt}{0pt}{MainBlue},
]
{\LARGE\bfseries\color{MainBlue} Elie ROUPHAEL, PhD}\\[3pt]
{\normalsize\color{DarkGray} Postdoctoral Researcher in Control, System Identification and Learning}\\[6pt]
{\footnotesize
\textcolor{AccentBlue}{\ding{41}}\ \href{mailto:elie.rouphael.98@gmail.com}{elie.rouphael.98@gmail.com}\quad
\textcolor{AccentBlue}{\textbullet}\ Grenoble, France\quad
\textcolor{AccentBlue}{\ding{226}}\ \href{https://elierouphael.github.io}{elierouphael.github.io}\\[2pt]
\textcolor{AccentBlue}{\ding{226}}\ \href{https://linkedin.com/in/elie-rouphael}{linkedin.com/in/elie-rouphael}\quad
\textcolor{AccentBlue}{\ding{226}}\ \href{https://github.com/ElieRouphael}{github.com/ElieRouphael}
}
\end{tcolorbox}

% ---------------------------------------------------------------- profile
\cvsection{Research profile}
\justifying
Postdoctoral researcher with a background in automatic control, mathematical modeling and data-driven analysis of dynamical systems. My doctoral work developed realization-theoretic and identification tools for stochastic switched and LPV state-space models, with emphasis on minimal representations, innovation forms and algorithms built from noisy data. My current postdoctoral work at GIPSA-lab, LIAS and Michelin concerns real-time estimation and model learning for vehicle dynamics and tire-road interaction. I am interested in rigorous methods that connect theory, algorithms and implementable estimation tools for complex dynamical systems.

% ---------------------------------------------------------------- education
\cvsection{Education}
\begin{tabularx}{\textwidth}{@{}L{2.45cm}Y@{}}
\cvitem{2021--2025}{\textbf{PhD in Automatic Control and Computer Science}, Universit\'e de Lille, CRIStAL, France. Thesis: \textit{Towards Stochastic Realization Theory for Linear Switched Models}. Advisors: Lotfi Belkoura and Mih\'aly Petreczky. Defended 24 October 2025.}
\cvitem{2020--2021}{\textbf{MSc in Automatic Control and Electrical Energy}, Universit\'e de Poitiers, LIAS, France. Ranked first in class. Research project on non-stationary signal detection and frequency estimation.}
\cvitem{2016--2021}{\textbf{Engineering Degree in Electrical and Industrial Control Systems}, Lebanese University, Faculty of Engineering, Lebanon.}
\end{tabularx}

% ---------------------------------------------------------------- experience
\cvsection{Professional and research experience}
\role{Postdoctoral Researcher - LabCom I-TireLab}{Nov 2025--present}{GIPSA-lab (Grenoble INP--UGA), LIAS (Univ.\ Poitiers) and Michelin, Grenoble}{
\begin{itemize}
  \item Real-time prediction of slippery road scenarios with hybrid modeling and machine learning, supervised by Guillaume Merc\`ere and John-Jairo Martinez-Molina.
  \item Physics-informed and data-driven models for tire-road interaction and vehicle dynamics; Kalman filtering combined with data-based parameter estimation for real-time use.
  \item Digital-twin-based methods for tire-road interaction modeling and dynamics prediction.
\end{itemize}}

\role{PhD Researcher}{2021--2025}{CRIStAL, Universit\'e de Lille, France}{
\begin{itemize}
  \item Stochastic realization results for generalized linear switched systems with inputs: deterministic and stochastic decompositions, innovation forms, minimality.
  \item Identification and realization algorithms for stochastic state-space models with switching signals, with numerical validation.
  \item Collaboration with SZTAKI, Budapest, on deep state-space and time-series models with generalization guarantees.
\end{itemize}}

\role{ATER and Teaching Assistant}{2022--2025}{Universit\'e de Lille, France}{
\begin{itemize}
  \item ATER (2024--2025): 195+ hours of teaching in Python, C, control theory, microcontrollers and GRAFCET; supervision of 3 student internships and 1 bachelor's project.
  \item Teaching assistant (2022--2024): 128+ hours of tutorials and labs in automatic control, signal processing and feedback systems.
\end{itemize}}

\role{Research Intern}{Mar--Jul 2021}{LIAS, Universit\'e de Poitiers, France}{
\begin{itemize}
  \item Detection of non-stationary harmonic disturbances in electrical systems (led to the IFAC 2023 paper below).
\end{itemize}}

% ---------------------------------------------------------------- publications
\newpage
\cvsection{Publications}
\begin{list}{}{%
  \setlength{\leftmargin}{2.1em}
  \setlength{\labelwidth}{1.9em}
  \setlength{\labelsep}{0.2em}
  \setlength{\itemsep}{3pt}
  \setlength{\topsep}{2pt}
  \setlength{\parsep}{0pt}
}
\item[\pubnum{1}] E. Rouphael, M. Mejari, M. Petreczky, L. Belkoura, ``Toward Stochastic Realization Theory for Generalized Linear Switched Systems With Inputs: Decomposition Into Stochastic and Deterministic Components and Existence and Uniqueness of Innovation Form,'' \textit{IEEE Control Systems Letters}, 2024. DOI: \href{https://doi.org/10.1109/LCSYS.2024.3417189}{10.1109/LCSYS.2024.3417189}.
\item[\pubnum{2}] E. Rouphael, M. Mejari, M. Petreczky, L. Belkoura, ``Minimal Covariance Realization and System Identification Algorithm for a Class of Stochastic Linear Switched Systems with i.i.d. Switching,'' \textit{IEEE Conference on Decision and Control}, 2024. DOI: \href{https://doi.org/10.1109/CDC56724.2024.10886432}{10.1109/CDC56724.2024.10886432}.
\item[\pubnum{3}] E. Rouphael, S. Cauet, A. Chamroo, ``Variable Frequency Monitoring of Power Grid: A Modified Observer Approach,'' \textit{IFAC-PapersOnLine}, 2023. DOI: \href{https://doi.org/10.1016/j.ifacol.2023.10.1830}{10.1016/j.ifacol.2023.10.1830}.
\item[\pubnum{4}] E. Rouphael, M. Petreczky, L. Belkoura, ``On Minimal LPV State-Space Representations in Innovation Form: An Algebraic Characterization,'' \textit{IEEE Conference on Decision and Control}, 2022. DOI: \href{https://doi.org/10.1109/CDC51059.2022.9992964}{10.1109/CDC51059.2022.9992964}.
\end{list}

% ---------------------------------------------------------------- projects
\cvsection{Supervised student projects}
\begin{itemize}
  \item \textbf{PAC bounds for stable state-space models:} generalization bounds for deep state-space models, with SZTAKI (Budapest).
  \item \textbf{Classical vs.\ deep vs.\ state-space time-series models:} benchmark of ARIMA/ARMAX, RNNs, CNNs and Transformers against SSMs such as Mamba.
  \item \textbf{Transport-mode prediction:} predicting individual transport choices from Hauts-de-France data for mobility planning.
  \item \textbf{System identification toolbox:} open-source Python tools for LPV, hybrid and switched systems.
\end{itemize}

% ---------------------------------------------------------------- skills
\cvsection{Technical skills}
\begin{tabularx}{\textwidth}{@{}L{3.1cm}Y@{}}
\cvitem{Programming}{\chip{Python} \chip{Matlab/Simulink} \chip{C/C++} \chip{SQL} \chip{R} \chip{Git}}
\cvitem{Scientific Python}{NumPy, SciPy, Pandas, Matplotlib, Scikit-learn; reproducible numerical experiments.}
\cvitem{Machine learning}{PyTorch, TensorFlow/Keras, JAX; neural networks, deep sequence models, physics-informed learning.}
\cvitem{Systems and control}{State-space modeling, stochastic systems, switched/LPV/hybrid models, observers, Kalman filtering, system identification, MPC.}
\end{tabularx}

% ---------------------------------------------------------------- certificates
\cvsection{Certificates and training}
\begin{tabularx}{\textwidth}{@{}L{3.1cm}Y@{}}
\cvitem{Universities}{\coursera{AI for Design and Optimization}{2D5GB3BJHFXB} (University of Michigan, 2026); \coursera{Bayesian Statistics: From Concept to Data Analysis}{I8P28MSPF616} (UC Santa Cruz, 2025); \coursera{Neural Networks and Deep Learning}{KYURTYL99Z7T} (DeepLearning.AI, 2023).}
\cvitem{IBM}{\coursera{Deep Learning and Reinforcement Learning}{GK0MYMKMM3KI}; \coursera{Specialized Models: Time Series and Survival Analysis}{0K0FGSEZE355}; \coursera{Deep Learning with Keras and TensorFlow}{NZJNT91GKH20}; \coursera{Introduction to Deep Learning \& Neural Networks with Keras}{TBF14TQEYMQV}; \coursera{Machine Learning with Python}{SSK9CKVWJ3BS}; \coursera{Generative AI and LLMs: Architecture and Data Preparation}{3Z6BM3POGK2Z}; \coursera{Python for Data Science, AI \& Development}{0XZF5ESFE0GM}; \coursera{Introduction to Data Engineering}{X209C9OT9Q9T}; \coursera{Python Project for Data Engineering}{K38NIHNO4DWG}; \coursera{Introduction to Relational Databases}{DIAYZGXY5KQL}; \coursera{Introduction to DevOps}{DF3EN2ZKCRY2}.}
\cvitem{Schools}{EECI course on Sparsity and Big Data in Control and ML; spring school on closed-loop and nonlinear system identification and deep learning; ACES summer school in modern control theory (60 h).}
\cvitem{Languages}{\chip{Arabic - native} \chip{French - B2, DELF} \chip{English - TOEIC 865/990}}
\end{tabularx}

% ---------------------------------------------------------------- references
\cvsection{References}
\noindent
\begin{tabularx}{\textwidth}{X X}
\textbf{Prof. Lotfi Belkoura} & \textbf{Dr. Mih\'aly Petreczky} \\
Professor, CRIStAL, Univ. Lille & CNRS Researcher, CRIStAL, Centrale Lille \\
\href{mailto:lotfi.belkoura@univ-lille.fr}{lotfi.belkoura@univ-lille.fr} &
\href{mailto:mihaly.petreczky@centralelille.fr}{mihaly.petreczky@centralelille.fr} \\
\\
\textbf{Prof. Guillaume Merc\`ere} & \textbf{Prof. John-Jairo Martinez-Molina} \\
Professor, LIAS, Univ. Poitiers & Professor, GIPSA-lab, Grenoble INP \\
\href{mailto:guillaume.mercere@univ-poitiers.fr}{guillaume.mercere@univ-poitiers.fr} &
\href{mailto:john-jairo.martinez-molina@grenoble-inp.fr}{john-jairo.martinez-molina@grenoble-inp.fr} \\
\end{tabularx}

\end{document}
